\section{Un teorema de existencia para los haces algebraicos coherentes}
\label{section:III.5}


\subsection{Enunciado del teorema}
\label{subsection:III.5.1}

\begin{env}[5.1.1]
\label{III.5.1.1}
Sea $A$ un anillo noetheriano ádico y sea $\mathfrak{J}$ un ideal de
definición de $A$, de modo que $A$ sea separado y completo para la
topología $\mathfrak{J}$-ádica. Si $Y=\Spec(A)$, el preesquema formal
afín $\Spf(A)$ se identifica con la compleción
$\widehat{Y}$ de $Y$ a lo largo del subconjunto cerrado
$Y'=\mathrm{V}(\mathfrak{J})$ \sref[I]{I.10.10.1}. Sea $X$ un
$Y$-preesquema ordinario de tipo finito, de morfismo estructural
$f:X\to Y$. Denotemos por $\widehat{X}$ la compleción de $X$ a lo largo
del subconjunto cerrado $X'=f^{-1}(Y')$, o, equivalentemente, el
$\widehat{Y}$-preesquema formal $X\times_Y\widehat{Y}$, y por
$\widehat{f}:\widehat{X}\to\widehat{Y}$ la extensión de $f$ a las
compleciones. Finalmente, para todo $\sh{O}_X$-módulo coherente
$\sh{F}$, denotemos por $\widehat{\sh{F}}$ su compleción
$\sh{F}_{/X'}$, que es un
$\sh{O}_{\widehat{X}}$-módulo coherente.
\end{env}

\begin{proposition}[5.1.2]
\label{III.5.1.2}
Con las hipótesis y la notación de \sref{III.5.1.1}, sea $\sh{F}$
un $\sh{O}_X$-módulo coherente cuyo soporte es propio sobre $Y$
\sref[II]{II.5.4.10}. Los homomorfismos canónicos
\sref{III.4.1.4}
\[
  \rho_i:
  \HH^i(X,\sh{F})
  \longrightarrow
  \HH^i(\widehat{X},\widehat{\sh{F}})
\]
son isomorfismos.
\end{proposition}

\begin{proof}
El $A$-módulo $\HH^i(X,\sh{F})$ es de tipo finito
\sref{III.3.2.4} y, por tanto, se identifica con su compleción
separada para la topología $\mathfrak{J}$-preádica
\sref[0]{0.7.3.6}. La proposición es, por consiguiente, un caso particular
de \sref{III.4.1.10}.
\end{proof}

Recordemos que los isomorfismos canónicos $\rho_i$ conmutan con los
homomorfismos de conexión asociados a toda sucesión exacta
\[
  0\longrightarrow\sh{F}'
  \longrightarrow\sh{F}
  \longrightarrow\sh{F}''
  \longrightarrow0
\]
de $\sh{O}_X$-módulos coherentes
\sref[0]{0.12.1.6} \sref[I]{I.10.8.9}.

\begin{corollary}[5.1.3]
\label{III.5.1.3}
Sean $\sh{F}$ y $\sh{G}$ $\sh{O}_X$-módulos coherentes tales que
la intersección de sus soportes sea propia sobre $Y$. El
homomorfismo canónico
\[
  \Hom_{\sh{O}_X}(\sh{F},\sh{G})
  \longrightarrow
  \Hom_{\sh{O}_{\widehat{X}}}
  (\widehat{\sh{F}},\widehat{\sh{G}})
  \tag{5.1.3.1}
  \label{III.5.1.3.1.eq}
\]
que envía un homomorfismo $u:\sh{F}\to\sh{G}$ a su compleción
$\widehat{u}:\widehat{\sh{F}}\to\widehat{\sh{G}}$, es un
isomorfismo. Además, cuando $f$ es cerrado,
$\widehat{u}$ es inyectivo (respectivamente, sobreyectivo) si y solo
si $u$ es inyectivo (respectivamente, sobreyectivo).
\end{corollary}

\begin{proof}
La primera afirmación es un caso particular de \sref{III.4.5.3}; también
se deduce de que la fuente de \eqref{III.5.1.3.1.eq} es
un $A$-módulo de tipo finito y, por tanto, se identifica con su
compleción separada. Para la segunda afirmación, observemos que, por
\sref[I]{I.10.8.14}, $\widehat{u}$ es inyectivo (respectivamente,
sobreyectivo) si y solo si existe un entorno de $X'$ en
$X$ sobre el cual $u$ es inyectivo (respectivamente, sobreyectivo). La
conclusión se deduce entonces del lema siguiente.
\end{proof}

\oldpage[III]{150}

\begin{lemma}[5.1.3.1]
\label{III.5.1.3.1}
Bajo las hipótesis de \sref{III.5.1.1}, supongamos además que
$f$ es cerrado. Todo entorno de $X'$ en $X$ es entonces igual
a $X$.
\end{lemma}

\begin{proof}
Nos reducimos primero al caso $f(X)=Y$. En efecto, por hipótesis,
$f(X)$ es un subconjunto cerrado $Z$ de $Y$. Podemos sustituir $f$ por
$f_{\mathrm{red}}$ \sref[I]{I.6.3.4} y suponer así que $X$
e $Y$ son reducidos; podemos entonces sustituir $Y$ por el subpreesquema
cerrado reducido cuyo espacio subyacente es $Z$
\sref[I]{I.5.2.2}, pues todo ideal de $A$ es cerrado y todo
anillo cociente de $A$ es, por tanto, ádico y noetheriano.

Se tiene entonces $f(X')=Y'$. Si $V$ es un entorno abierto de $X'$
en $X$, el conjunto $f(X-V)$ es cerrado en $Y$ y no corta a $Y'$.
Esto es imposible salvo que $X-V$ sea vacío, pues $\mathfrak{J}$
está contenido en el radical de $A$
\sref[I]{I.1.1.15} \sref[0]{0.7.1.10}. Por tanto, $V=X$.
\end{proof}

Si nos restringimos a los $\sh{O}_X$-módulos coherentes cuyo
soporte es propio sobre $Y$, \sref{III.5.1.3} afirma, en lenguaje
categórico, que el funtor
\[
  \sh{F}\longmapsto\widehat{\sh{F}}
\]
es plenamente fiel de la categoría de tales $\sh{O}_X$-módulos en
la categoría de los $\sh{O}_{\widehat{X}}$-módulos coherentes. Proporciona así
una equivalencia de la primera categoría con una subcategoría plena
de la segunda \sref[0]{0.8.1.6}. El teorema de existencia
mostrará que, cuando $X$ es propio sobre $Y$, esta subcategoría es en
realidad la categoría de todos los
$\sh{O}_{\widehat{X}}$-módulos coherentes. Más precisamente:

\begin{theorem}[5.1.4]
\label{III.5.1.4}
Sea $A$ un anillo noetheriano ádico, sea $Y=\Spec(A)$, sea
$\mathfrak{J}$ un ideal de definición de $A$, y pongamos
$Y'=\mathrm{V}(\mathfrak{J})$. Sea $f:X\to Y$ un morfismo separado
de tipo finito y pongamos $X'=f^{-1}(Y')$. Sean
\[
  \widehat{Y}=Y_{/Y'}=\Spf(A),\qquad
  \widehat{X}=X_{/X'}
\]
las compleciones de $Y$ y $X$ a lo largo de $Y'$ y $X'$, y sea
$\widehat{f}:\widehat{X}\to\widehat{Y}$ la extensión de $f$ a
las compleciones. Entonces el funtor
\[
  \sh{F}\longmapsto\sh{F}_{/X'}=\widehat{\sh{F}}
\]
es una equivalencia de la categoría de los
$\sh{O}_X$-módulos coherentes cuyo soporte es propio sobre $\Spec(A)$ en la
categoría de los $\sh{O}_{\widehat{X}}$-módulos coherentes cuyo soporte
es propio sobre $\Spf(A)$.
\end{theorem}

En otras palabras, teniendo en cuenta \sref{III.5.1.3}:

\begin{corollary}[5.1.5]
\label{III.5.1.5}
Un $\sh{O}_{\widehat{X}}$-módulo es isomorfo a la compleción de
un $\sh{O}_X$-módulo coherente cuyo soporte es propio sobre
$\Spec(A)$ si y solo si es coherente y su soporte es propio
sobre $\Spf(A)$.
\end{corollary}

El caso más importante es el siguiente:

\begin{corollary}[5.1.6]
\label{III.5.1.6}
Supongamos que $X$ es propio sobre $Y=\Spec(A)$. Entonces el funtor
$\sh{F}\mapsto\widehat{\sh{F}}$ es una equivalencia entre la
categoría de los $\sh{O}_X$-módulos coherentes y la categoría de los
$\sh{O}_{\widehat{X}}$-módulos coherentes.
\end{corollary}

\begin{env}[5.1.7]
\label{III.5.1.7}
\emph{Escolio.}
Teniendo en cuenta la caracterización de los haces coherentes sobre
preesquemas formales \sref[I]{I.10.11.3}, vemos que, bajo las
condiciones de \sref{III.5.1.1}, dar un
$\sh{O}_X$-módulo coherente cuyo soporte es propio sobre $\Spec(A)$
equivale---si
\[
  Y_n=\Spec(A/\mathfrak{J}^{n+1}),\qquad
  X_n=X\times_Y Y_n
\]
---a dar un sistema proyectivo $(\sh{F}_n)$ de
$\sh{O}_{X_n}$-módulos coherentes tal que, para $m\leq n$,
\[
  \sh{F}_m
  =
  \sh{F}_n\otimes_{\sh{O}_{Y_n}}\sh{O}_{Y_m},
  \qquad\text{o equivalentemente}\qquad
  \sh{F}_m
  =
  \sh{F}_n/\mathfrak{J}^{m+1}\sh{F}_n,
\]
y tal que el soporte de $\sh{F}_0$ sea un subconjunto de $X_0$
propio sobre $Y_0$. Por \sref[I]{I.10.11.4}, los homomorfismos de
$\sh{O}_X$-módulos coherentes se interpretan de modo análogo como
homomorfismos de sistemas proyectivos de
$\sh{O}_{X_n}$-módulos coherentes.

\oldpage[III]{151}

En todas las aplicaciones conocidas, $A$ es de hecho un anillo
noetheriano ádico local. Así, los $Y_n$ son espectros de anillos
artinianos locales, y los resultados de esta subsección y de las
anteriores reducen, en gran medida, la geometría algebraica sobre un
anillo noetheriano ádico local a la geometría algebraica sobre anillos
artinianos locales.
\end{env}

\begin{corollary}[5.1.8]
\label{III.5.1.8}
Bajo las hipótesis de \sref{III.5.1.4}, la aplicación
\[
  Z\longmapsto\widehat{Z}=Z_{/(Z\cap X')}
\]
es una biyección del conjunto de los subpreesquemas cerrados $Z$ de $X$
que son propios sobre $Y$ sobre el conjunto de los subpreesquemas
formales cerrados de $\widehat{X}$ que son propios sobre
$\widehat{Y}$.
\end{corollary}

\begin{proof}
Un subpreesquema formal cerrado de $\widehat{X}$ es de la forma
\[
  \bigl(T,
    (\sh{O}_{\widehat{X}}/\sh{A})|_T\bigr),
\]
donde $\sh{A}$ es un ideal coherente de
$\sh{O}_{\widehat{X}}$ \sref[I]{I.10.14.2}. Si $T$ es propio sobre
$\widehat{Y}$, \sref{III.5.1.4} muestra que
$\sh{O}_{\widehat{X}}/\sh{A}$ es isomorfo a un
$\sh{O}_{\widehat{X}}$-módulo de la forma
$\widehat{\sh{F}}$, donde $\sh{F}$ es un
$\sh{O}_X$-módulo coherente cuyo soporte es propio sobre $Y$. Además,
\sref{III.5.1.3} muestra que el homomorfismo canónico
\[
  \sh{O}_{\widehat{X}}
  \longrightarrow
  \sh{O}_{\widehat{X}}/\sh{A}
\]
es de la forma $\widehat{u}$, donde
$u:\sh{O}_X\to\sh{F}$ es un homomorfismo sobreyectivo de
$\sh{O}_X$-módulos. Por consiguiente,
$\sh{F}=\sh{O}_X/\sh{N}$ para un ideal coherente $\sh{N}$ de
$\sh{O}_X$, y
$\sh{A}=\widehat{\sh{N}}$ \sref[I]{I.10.8.8}. La conclusión se
deduce ahora de \sref[I]{I.10.14.7}.
\end{proof}


\subsection{Demostración del teorema de existencia: caso proyectivo y cuasiproyectivo}
\label{subsection:III.5.2}

\begin{env}[5.2.1]
\label{III.5.2.1}
Bajo las condiciones de \sref{III.5.1.4}, diremos
provisionalmente que un $\sh{O}_{\widehat{X}}$-módulo coherente es
\emph{algebraizable} si es isomorfo a la compleción
$\widehat{\sh{F}}$ de un $\sh{O}_X$-módulo coherente $\sh{F}$ cuyo
soporte es propio sobre $Y$.
\end{env}

\begin{lemma}[5.2.2]
\label{III.5.2.2}
Sean $\sh{F}'$ y $\sh{G}'$ $\sh{O}_{\widehat{X}}$-módulos
algebraizables. Para todo homomorfismo
$u:\sh{F}'\to\sh{G}'$, los módulos
$\Ker(u)$, $\Im(u)$ y $\Coker(u)$ son algebraizables.
\end{lemma}

\begin{proof}
Escribamos
\[
  \sh{F}'=\widehat{\sh{F}},\qquad
  \sh{G}'=\widehat{\sh{G}},
\]
donde $\sh{F}$ y $\sh{G}$ son $\sh{O}_X$-módulos coherentes cuyos
soportes son propios sobre $Y$. Por \sref{III.5.1.3},
$u=\widehat{v}$ para un homomorfismo
$v:\sh{F}\to\sh{G}$. Como la compleción es exacta,
$\Ker(\widehat{v})\simeq\widehat{\Ker(v)}$; y como el soporte
de $\Ker(v)$ está contenido en el de $\sh{F}$, se deduce que
$\Ker(u)$ es algebraizable. Los argumentos para $\Im(u)$ y
$\Coker(u)$ son los mismos.
\end{proof}

\begin{proposition}[5.2.3]
\label{III.5.2.3}
Sea $A$ un anillo noetheriano ádico, sea $\mathfrak{J}$ un ideal de
definición de $A$, sea $\mathfrak{Y}=\Spf(A)$, y sea
$f:\mathfrak{X}\to\mathfrak{Y}$ un morfismo propio de preesquemas
formales. Pongamos
\[
  Y_k=\Spec(A/\mathfrak{J}^{k+1}),\qquad
  X_k=\mathfrak{X}\times_{\mathfrak{Y}}Y_k,
\]
y, para todo $\sh{O}_{\mathfrak{X}}$-módulo $\sh{F}$, pongamos
\[
  \sh{F}_k
  =
  \sh{F}\otimes_{\sh{O}_{\mathfrak{X}}}\sh{O}_{X_k}
  =
  \sh{F}/\mathfrak{J}^{k+1}\sh{F}.
\]
Sea $\sh{L}$ un $\sh{O}_{\mathfrak{X}}$-módulo invertible, y
supongamos que
$\sh{L}_0=\sh{L}/\mathfrak{J}\sh{L}$ es un
$\sh{O}_{X_0}$-módulo amplio. Para todo
$\sh{O}_{\mathfrak{X}}$-módulo $\sh{F}$ y todo entero $n$, pongamos
\[
  \sh{F}(n)=\sh{F}\otimes\sh{L}^{\otimes n}.
\]
Entonces, para todo $\sh{O}_{\mathfrak{X}}$-módulo coherente $\sh{F}$,
existe un entero $n_0$ tal que, para todo $n\geq n_0$:
\begin{enumerate}
  \item[\rm{(i)}] el homomorfismo canónico
    \[
      \HH^0(\mathfrak{X},\sh{F}(n))
      \longrightarrow
      \HH^0(\mathfrak{X},\sh{F}_k(n))
    \]
    es sobreyectivo para todo $k\geq0$;
  \item[\rm{(ii)}]
    $\HH^q(\mathfrak{X},\sh{F}(n))=0$ para todo $q>0$.
\end{enumerate}
\end{proposition}

\begin{proof}
Los espacios subyacentes a $\mathfrak{X}$ y $X_0$ son los mismos.
Los haces
\[
  \sh{M}_k
  =
  \mathfrak{J}^k\sh{F}/\mathfrak{J}^{k+1}\sh{F},
\]
al ser anulados por $\mathfrak{J}$, pueden considerarse
$\sh{O}_{X_0}$-módulos coherentes \sref[0]{0.5.3.10}. Además, si
\[
  \sh{M}_k(n)
  =
  \sh{M}_k\otimes_{\sh{O}_{X_0}}\sh{L}_0^{\otimes n},
\]
entonces
\[
  \sh{M}_k(n)
  =
  \mathfrak{J}^k\sh{F}(n)/
  \mathfrak{J}^{k+1}\sh{F}(n).
\]
Como $\sh{L}_0$ es amplio para $f_0:X_0\to Y_0$ y $f_0$ es
propio,
\oldpage[III]{152}
$f_0$ es proyectivo \sref[II]{II.5.5.4}. Sea
\[
  S=\bigoplus_{k\geq0}
  \mathfrak{J}^k/\mathfrak{J}^{k+1}
\]
la $A$-álgebra graduada asociada a la
filtración $\mathfrak{J}$-ádica de $A$. Es de tipo finito
porque $A$ es noetheriano. Pongamos
\[
  \sh{S}'=f_0^*(\widetilde{S}),\qquad
  \sh{M}=\bigoplus_{k\geq0}\sh{M}_k.
\]
Entonces $\sh{S}'$ es una $\sh{O}_{X_0}$-álgebra cuasi-coherente de
tipo finito, y $\sh{M}$ es un $\sh{S}'$-módulo graduado
cuasi-coherente de tipo finito, pues $\sh{F}_0$ es coherente y
genera $\sh{M}$ como $\sh{S}'$-módulo. El Teorema
\sref{III.2.4.1},~\emph{(ii)}, proporciona, por tanto, un entero $n_0$
tal que, para $n\geq n_0$ y todo $k$,
\[
  \HH^q(X_0,\sh{M}_k(n))=0
  \qquad(q>0).
  \tag{5.2.3.1}
  \label{III.5.2.3.1}
\]
Considerado ahora como $\sh{O}_{\mathfrak{X}}$-módulo,
$\sh{M}_k(n)$ satisface, por consiguiente, también
\[
  \HH^q(\mathfrak{X},\sh{M}_k(n))=0
  \qquad(q>0,\ n\geq n_0).
\]

Apliquemos la sucesión exacta de cohomología a
\[
  0\longrightarrow
  \frac{\mathfrak{J}^k\sh{F}(n)}
       {\mathfrak{J}^{k+1}\sh{F}(n)}
  \longrightarrow
  \frac{\mathfrak{J}^h\sh{F}(n)}
       {\mathfrak{J}^{k+1}\sh{F}(n)}
  \longrightarrow
  \frac{\mathfrak{J}^h\sh{F}(n)}
       {\mathfrak{J}^{k}\sh{F}(n)}
  \longrightarrow0.
\]
Para $0\leq h<k$, $n\geq n_0$ y $q>0$, una inducción sobre $k-h$
da
\[
  \HH^q\left(
    \mathfrak{X},
    \frac{\mathfrak{J}^h\sh{F}(n)}
         {\mathfrak{J}^{k}\sh{F}(n)}
  \right)=0,
  \tag{5.2.3.2}
  \label{III.5.2.3.2}
\]
y, en particular, para $h=0$,
\[
  \HH^q(\mathfrak{X},\sh{F}_k(n))=0
  \qquad(n\geq n_0,\ k\geq0,\ q>0).
  \tag{5.2.3.3}
  \label{III.5.2.3.3}
\]
Otra parte de la sucesión exacta de cohomología, con $h=0$, da
\[
  \HH^0(\mathfrak{X},\sh{F}_{k+1}(n))
  \longrightarrow
  \HH^0(\mathfrak{X},\sh{F}_k(n))
  \longrightarrow
  \HH^1\left(
    \mathfrak{X},
    \frac{\mathfrak{J}^k\sh{F}(n)}
         {\mathfrak{J}^{k+1}\sh{F}(n)}
  \right)
  =0.
  \tag{5.2.3.4}
  \label{III.5.2.3.4}
\]
Se deduce que, para $h\leq k$, la aplicación canónica
\[
  \HH^0(\mathfrak{X},\sh{F}_k(n))
  \longrightarrow
  \HH^0(\mathfrak{X},\sh{F}_h(n))
  \tag{5.2.3.5}
  \label{III.5.2.3.5}
\]
es sobreyectiva. Así, para todo $q$, el sistema proyectivo
\[
  \bigl(\HH^q(\mathfrak{X},\sh{F}_k(n))\bigr)_{k\geq0}
\]
satisface la condición (ML) cuando $n\geq n_0$.

Además, todo abierto formal afín $U$ de $\mathfrak{X}$ es un
abierto afín de cada $X_k$ \sref[I]{I.10.5.2}. Por tanto,
$\HH^q(U,\sh{F}_k(n))=0$ para $q>0$ \sref[I]{I.1.3.1}, y
\[
  \HH^0(U,\sh{F}_k(n))
  \longrightarrow
  \HH^0(U,\sh{F}_h(n))
\]
es sobreyectivo para $h\leq k$ \sref[I]{I.1.3.9}. Por tanto, se satisfacen
las hipótesis de \sref[0]{0.13.3.1}. Para $n\geq n_0$
obtenemos:
\begin{enumerate}
  \item[$1^\circ$] para todo $q>0$, la aplicación
    \[
      \HH^q(\mathfrak{X},\sh{F}(n))
      \longrightarrow
      \varprojlim_k
      \HH^q(\mathfrak{X},\sh{F}_k(n))
    \]
    es biyectiva; por tanto, por \eqref{III.5.2.3.3},
    $\HH^q(\mathfrak{X},\sh{F}(n))=0$;
  \item[$2^\circ$] el homomorfismo
    \[
      \HH^0(\mathfrak{X},\sh{F}(n))
      \longrightarrow
      \varprojlim_k
      \HH^0(\mathfrak{X},\sh{F}_k(n))
    \]
    es biyectivo. Como los homomorfismos
    \eqref{III.5.2.3.5} son sobreyectivos, también lo es cada homomorfismo
    \[
      \varprojlim_k
      \HH^0(\mathfrak{X},\sh{F}_k(n))
      \longrightarrow
      \HH^0(\mathfrak{X},\sh{F}_h(n)).
    \]
\end{enumerate}
Esto demuestra las dos afirmaciones.
\end{proof}

\begin{corollary}[5.2.4]
\label{III.5.2.4}
Bajo las hipótesis de \sref{III.5.2.3}, para todo
$\sh{O}_{\mathfrak{X}}$-módulo coherente $\sh{F}$ existe un entero $N$
tal que, para $n\geq N$, $\sh{F}(n)$ está generado por sus secciones
sobre
\oldpage[III]{153}
$\mathfrak{X}$. Equivalentemente, $\sh{F}$ es isomorfo a un cociente
de un $\sh{O}_{\mathfrak{X}}$-módulo de la forma
$(\sh{O}_{\mathfrak{X}}(-n))^k$.
\end{corollary}

\begin{proof}
Como $X_0$ es noetheriano, la hipótesis sobre $\sh{L}_0$ y
\sref[II]{II.4.5.5} proporcionan un entero $n_0$ tal que
$\sh{F}_0(n)$ está generado por sus secciones sobre $\mathfrak{X}$ para
$n\geq n_0$. Podemos tomar $n_0$ suficientemente grande para que
\[
  \Gamma(\mathfrak{X},\sh{F}(n))
  \longrightarrow
  \Gamma(\mathfrak{X},\sh{F}_0(n))
\]
sea también sobreyectivo para $n\geq n_0$ \sref{III.5.2.3}. Existen,
por tanto, un número finito de secciones
$s_i\in\Gamma(\mathfrak{X},\sh{F}(n))$ cuyas imágenes generan
$\sh{F}_0(n)$ \sref[0]{0.5.2.3}. Como $\mathfrak{J}$ está
contenido en el ideal maximal del anillo local en todo punto
de $\mathfrak{X}$, el lema de Nakayama aplicado a esos anillos locales
muestra que los $s_i$ generan $\sh{F}(n)$
\sref[0]{0.5.1.1}.
\end{proof}

\begin{env}[5.2.5]
\label{III.5.2.5}
\emph{Demostración del teorema de existencia: caso proyectivo.}
Con la notación de \sref{III.5.1.4}, supongamos que $f$ es
proyectivo, de modo que existe un $\sh{O}_X$-módulo amplio $\sh{L}$
\sref[I]{I.5.5.4}. Por definición,
\[
  X_n=\widehat{X}\times_{\widehat{Y}}Y_n
\]
es el subpreesquema cerrado
$X\times_Y Y_n=f^{-1}(Y_n)$ de $X$. Si
\[
  \sh{L}'=\sh{L}\otimes_{\sh{O}_X}\sh{O}_{\widehat{X}}
\]
es la compleción de $\sh{L}$, entonces
$\sh{L}'_0=\sh{L}/\mathfrak{J}\sh{L}$, considerado como
$\sh{O}_{X_0}$-módulo; este módulo es amplio
\sref[II]{II.4.6.13}[(i~\emph{bis})].

Podemos, por tanto, aplicar \sref{III.5.2.4} a $\sh{L}'$ y a todo
$\sh{O}_{\widehat{X}}$-módulo coherente $\sh{F}$. Se deduce que
$\sh{F}$ es isomorfo a un cociente de
\[
  \sh{G}=
  \bigl(\sh{L}'^{\otimes(-n)}\bigr)^k
\]
para enteros adecuados $n>0$ y $k>0$. El módulo $\sh{G}$ es la
compleción de $(\sh{L}^{\otimes(-n)})^k$
\sref[I]{I.10.8.10} y, por tanto, es algebraizable. Sea
$u:\sh{G}\to\sh{F}$ la sobreyección canónica, y pongamos
$\sh{H}=\Ker(u)$; este es un
$\sh{O}_{\widehat{X}}$-módulo coherente \sref[0]{0.5.3.4}. Aplicando el mismo
argumento a $\sh{H}$, obtenemos un
$\sh{O}_{\widehat{X}}$-módulo algebraizable $\sh{K}$ y un homomorfismo
$v:\sh{K}\to\sh{G}$ tal que $\sh{H}=\Im(v)$. Así,
$\sh{F}=\Coker(v)$, y $\sh{F}$ es algebraizable por
\sref{III.5.2.2}.
\end{env}

\begin{env}[5.2.6]
\label{III.5.2.6}
\emph{Demostración del teorema de existencia: caso cuasiproyectivo.}
Conservemos la notación de \sref{III.5.1.4} y supongamos ahora que $f$
es cuasiproyectivo. Existe un morfismo proyectivo
$g:Z\to Y$ tal que $X$ se identifica con el $Y$-preesquema
inducido sobre un subconjunto abierto de $Z$ \sref[II]{II.5.3.2}. Si
$Z'=g^{-1}(Y')$, entonces $X'=X\cap Z'$. Por consiguiente, la compleción
$\widehat{X}=X_{/X'}$ se identifica con el preesquema formal
inducido por $\widehat{Z}=Z_{/Z'}$ sobre el subconjunto abierto $X\cap Z'$ de
$\widehat{Z}$ \sref[I]{I.10.8.5}.

Sea $\sh{F}'$ un $\sh{O}_{\widehat{X}}$-módulo coherente cuyo
soporte $T'$ sea propio sobre $\widehat{Y}$. Por definición, existe
un subpreesquema cerrado de $X'$ de espacio subyacente
$T'\subset X'$ tal que la restricción $T'\to Y'$ de $f$ es
propia. Por tanto, $T'$ es propio sobre $Y$ y, en consecuencia, cerrado en
$Z'$ \sref[II]{II.5.4.10}. Se deduce que $\sh{F}'$ es inducido
sobre $\widehat{X}$ por el $\sh{O}_{\widehat{Z}}$-módulo $\sh{G}'$
obtenido pegando $\sh{F}'$, sobre el abierto $\widehat{X}$, con el
haz cero sobre el abierto $\widehat{Z}-T'$; ambos haces coinciden sobre
su intersección $\widehat{X}-T'$.

El módulo $\sh{G}'$ es coherente. Por \sref{III.5.2.5}, existe un
$\sh{O}_Z$-módulo coherente $\sh{G}$ tal que
$\sh{G}'=\widehat{\sh{G}}$. Sea $T$ el soporte de $\sh{G}$;
entonces $T'=T\cap Z'$ \sref[I]{I.10.8.12}. Si $h$ es la restricción
de $g$ al subpreesquema cerrado reducido de $Z$ cuyo espacio subyacente
es $T$, entonces
\[
  T'=h^{-1}(Y')=T\cap g^{-1}(Y').
\]
Así, $X\cap T$ es un subconjunto abierto de $T$ que contiene a $T'$.

\oldpage[III]{154}
Como $h$ es propio \sref[II]{II.5.4.2}, y por tanto cerrado,
\sref{III.5.1.3.1} da $X\cap T=T$. Así, $T\subset X$; y
como $T$ es cerrado en $Z$, es propio sobre $Y$. Si $\sh{F}$ es
el haz inducido sobre $X$ por $\sh{G}$, su compleción
$\widehat{\sh{F}}$ es inducida sobre $\widehat{X}$ por
$\widehat{\sh{G}}$ \sref[I]{I.10.8.4} y, por tanto, es igual a
$\sh{F}'$. Esto termina la demostración.
\end{env}


\subsection{Demostración del teorema de existencia: caso general}
\label{subsection:III.5.3}

\begin{lemma}[5.3.1]
\label{III.5.3.1}
Bajo las condiciones de \sref{III.5.1.4}, sea
\[
  0\longrightarrow\sh{H}\longrightarrow\sh{F}
  \longrightarrow\sh{G}\longrightarrow0
\]
una sucesión exacta de $\sh{O}_{\widehat{X}}$-módulos coherentes.
Si $\sh{G}$ y $\sh{H}$ son algebraizables, entonces $\sh{F}$ es
algebraizable.
\end{lemma}

\begin{proof}
Supongamos que
\[
  \sh{G}=\widehat{\sh{B}},
  \qquad
  \sh{H}=\widehat{\sh{C}},
\]
donde $\sh{B}$ y $\sh{C}$ son $\sh{O}_X$-módulos coherentes cuyos
soportes son propios sobre $Y$. El módulo $\sh{F}$ define canónicamente
un elemento del $A$-módulo
\[
  \Ext^1_{\sh{O}_{\widehat{X}}}
  (\widehat{X};\widehat{\sh{B}},\widehat{\sh{C}})
\]
\sref[0]{0.12.3.2}. Las hipótesis implican que este $A$-módulo es
canónicamente isomorfo a
\[
  \Ext^1_{\sh{O}_X}(X;\sh{B},\sh{C})
\]
\sref{III.4.5.2}. Existe, por tanto, una sucesión exacta
\[
  0\longrightarrow\sh{C}\longrightarrow\sh{A}
  \longrightarrow\sh{B}\longrightarrow0
\]
de $\sh{O}_X$-módulos coherentes cuya clase canónica se envía en la
clase correspondiente a $\sh{F}$. Por definición, teniendo en cuenta
\sref[I]{I.10.8.8}, (ii), esto significa que $\sh{F}$ es isomorfo a
$\widehat{\sh{A}}$. El lema se deduce, pues
$\Supp(\sh{A})$ está contenido en
$\Supp(\sh{B})\cup\Supp(\sh{C})$ y es, por tanto, propio sobre $Y$.
\end{proof}

\begin{corollary}[5.3.2]
\label{III.5.3.2}
Bajo las condiciones de \sref{III.5.1.1}, sea
$u:\sh{F}\to\sh{G}$ un homomorfismo de
$\sh{O}_{\widehat{X}}$-módulos coherentes. Si $\sh{G}$, $\Ker(u)$ y
$\Coker(u)$ son algebraizables, entonces $\sh{F}$ es algebraizable.
\end{corollary}

\begin{proof}
El Lema \sref{III.5.2.2}, aplicado a
$\sh{G}\to\Coker(u)$, muestra que $\Im(u)$ es algebraizable. Basta entonces
aplicar el Lema \sref{III.5.3.1} a la sucesión exacta
\[
  0\longrightarrow\Ker(u)\longrightarrow\sh{F}
  \longrightarrow\Im(u)\longrightarrow0.
\]
\end{proof}

\begin{lemma}[5.3.3]
\label{III.5.3.3}
Bajo las condiciones de \sref{III.5.1.1}, sea $h:Z\to Y$ un
morfismo de tipo finito, sea $\widehat{Z}$ la compleción de $Z$
a lo largo de $Z'=h^{-1}(Y')$, sea $g:Z\to X$ un $Y$-morfismo propio,
y sea $\widehat{g}:\widehat{Z}\to\widehat{X}$ su extensión a
las compleciones. Para todo $\sh{O}_{\widehat{Z}}$-módulo
algebraizable $\sh{F}'$, el módulo
$\widehat{g}_*(\sh{F}')$ es un
$\sh{O}_{\widehat{X}}$-módulo algebraizable.
\end{lemma}

\begin{proof}
Si $\sh{F}$ es un $\sh{O}_Z$-módulo coherente tal que
$\sh{F}'=\widehat{\sh{F}}$, el primer teorema de comparación
\sref{III.4.1.5} identifica
$\widehat{g}_*(\widehat{\sh{F}})$ con la compleción de
$g_*(\sh{F})$.
\end{proof}

\begin{lemma}[5.3.4]
\label{III.5.3.4}
Sea $X$ un esquema ordinario noetheriano, sea $X'$ un subconjunto
cerrado de $X$, sea $f:Z\to X$ un morfismo propio, y pongamos
\[
  Z'=f^{-1}(X'),\qquad
  \widehat{X}=X_{/X'},\qquad
  \widehat{Z}=Z_{/Z'}.
\]
Sea $\widehat{f}:\widehat{Z}\to\widehat{X}$ la extensión de $f$
a las compleciones. Sea $\sh{M}$ un ideal coherente de
$\sh{O}_X$ tal que, si
\[
  U=X-\Supp(\sh{O}_X/\sh{M}),
\]
la restricción $f^{-1}(U)\to U$ de $f$ es un isomorfismo. Entonces,
para todo $\sh{O}_{\widehat{X}}$-módulo coherente $\sh{F}$, existe
un entero $n>0$ tal que el núcleo y el conúcleo del homomorfismo
canónico
\[
  \rho_{\sh{F}}:
  \sh{F}\longrightarrow
  \widehat{f}_*\bigl(\widehat{f}^{\,*}(\sh{F})\bigr)
\]
\sref[0]{0.4.4.3} son anulados por
$\widehat{\sh{M}}^{\,n}$.
\end{lemma}

\begin{proof}
Podemos restringirnos al caso $X=\Spec(B)$, donde $B$ es un
anillo noetheriano; entonces $X'=\mathrm{V}(\mathfrak{K})$ para un ideal
$\mathfrak{K}$ de $B$. Nos reducimos primero al caso en que $B$ es
noetheriano ádico y $\mathfrak{K}$ es un ideal de definición.

Sea $B_1$ la compleción separada de $B$ para la
topología $\mathfrak{K}$-preádica.
\oldpage[III]{155}
Si $\mathfrak{K}_1=\mathfrak{K}B_1$, entonces $B_1$ es un anillo
noetheriano ádico y $\mathfrak{K}_1$ es un ideal de definición.
Pongamos $X_1=\Spec(B_1)$, y sea $h:X_1\to X$ el morfismo
correspondiente al homomorfismo canónico $B\to B_1$. Si
$X'_1=h^{-1}(X')$, entonces $X'_1=\mathrm{V}(\mathfrak{K}_1)$. Finalmente,
pongamos
\[
  Z_1=Z\times_X X_1,\qquad f_1:Z_1\longrightarrow X_1.
\]
El morfismo $f_1$ es propio \sref[II]{II.5.4.2}. Denotemos por
$\widehat{X}_1$ la compleción de $X_1$ a lo largo de $X'_1$, por
\[
  \widehat{Z}_1=
  Z_1\times_{X_1}\widehat{X}_1
\]
la compleción de $Z_1$ a lo largo de
$Z'_1=f_1^{-1}(X'_1)$, y por
$\widehat{f}_1$ la extensión de $f_1$ a las compleciones.

La extensión
$\widehat{h}:\widehat{X}_1\to\widehat{X}$ corresponde a la
aplicación identidad de $B_1$ y es, por tanto, un isomorfismo
\sref[I]{I.10.9.1}. El morfismo correspondiente
$\widehat{Z}_1\to\widehat{Z}$ es también, por consiguiente, un isomorfismo,
y estos isomorfismos identifican $\widehat{f}_1$ con
$\widehat{f}$. Además,
$\sh{M}_1=h^*(\sh{M})$ es un ideal coherente de $\sh{O}_{X_1}$ y
\[
  \Supp(\sh{O}_{X_1}/\sh{M}_1)
  =h^{-1}\bigl(\Supp(\sh{O}_X/\sh{M})\bigr)
\]
\sref[I]{I.9.1.13}. Así, si
$U_1=X_1-\Supp(\sh{O}_{X_1}/\sh{M}_1)$, entonces
$U_1=h^{-1}(U)$, y la restricción
$f_1^{-1}(U_1)\to U_1$ es un isomorfismo
\sref[I]{I.3.2.7}. Finalmente, $\widehat{\sh{M}}$ y
$\widehat{\sh{M}}_1$ son identificados por $\widehat{h}$
\sref[I]{I.10.9.5}. Todas las hipótesis del lema son, pues,
satisfechas por $X_1$, $X'_1$, $f_1$ y $\sh{M}_1$; en lo sucesivo
podemos suponer que $B$ es noetheriano ádico y que $\mathfrak{K}$ es
un ideal de definición.

Se tiene entonces $\widehat{X}=\Spf(B)$ y
$\sh{F}=N^\Delta$, donde $N$ es un $B$-módulo de tipo finito. Así,
$\sh{F}=\widehat{\sh{G}}$, donde $\sh{G}$ es el
$\sh{O}_X$-módulo coherente $\widetilde{N}$ \sref[I]{I.10.10.5}, y
\[
  \widehat{f}^{\,*}(\sh{F})
  =\widehat{f^*(\sh{G})}
\]
\sref[I]{I.10.9.5}. Por el primer teorema de comparación
\sref{III.4.1.5},
\[
  \widehat{f}_*
  \bigl(\widehat{f^*(\sh{G})}\bigr)
  \simeq
  \widehat{f_*\bigl(f^*(\sh{G})\bigr)},
\]
y, por \sref{III.5.1.3}, el homomorfismo canónico
$\rho_{\sh{F}}$ es precisamente $\widehat{\rho}_{\sh{G}}$.

El núcleo $\sh{P}$ y el conúcleo $\sh{R}$ de
\[
  \rho_{\sh{G}}:
  \sh{G}\longrightarrow f_*\bigl(f^*(\sh{G})\bigr)
\]
son $\sh{O}_X$-módulos coherentes, y sus restricciones a $U$ son
nulas. Existe, por tanto, un entero $n>0$ tal que
\[
  \sh{M}^n\sh{P}=\sh{M}^n\sh{R}=0
\]
\sref[I]{I.9.3.4}. Se deduce que
\[
  \widehat{\sh{M}}^{\,n}\widehat{\sh{P}}
  =
  \widehat{\sh{M}}^{\,n}\widehat{\sh{R}}
  =0
\]
\sref[I]{I.10.8.8} \sref[I]{I.10.8.10}, lo que demuestra el lema.
\end{proof}

\begin{env}[5.3.5]
\label{III.5.3.5}
\emph{Fin de la demostración del teorema de existencia.}
Conservemos las hipótesis de \sref{III.5.1.4}. Utilicemos la inducción
noetheriana \sref[0]{0.2.2.2}, suponiendo conocido el teorema para todo
subpreesquema cerrado $T$ de $X$ cuyo espacio subyacente es distinto
de $X$; aquí $\widehat{T}$ denota la compleción de $T$ a lo largo de
$T\cap X'$. Podemos suponer que $X$ no es vacío.

Como $f$ es separado y de tipo finito, el lema de Chow
\sref[II]{II.5.6.1} proporciona un $Y$-esquema $Z$ y un $Y$-morfismo
$g:Z\to X$ tales que el morfismo estructural $h:Z\to Y$ es
cuasiproyectivo, $g$ es proyectivo y sobreyectivo, y existe un
subconjunto abierto no vacío $U$ de $X$ para el cual
$g^{-1}(U)\to U$ es un isomorfismo. Sea $\sh{M}$ un ideal coherente
de $\sh{O}_X$ que define un subpreesquema cerrado de espacio subyacente
$X-U$ \sref[I]{I.5.2.2}, y sea $\sh{F}$ un
$\sh{O}_{\widehat{X}}$-módulo coherente cuyo soporte $E$ es propio sobre $Y$.
Denotemos por $\widehat{Z}$ la compleción de $Z$ a lo largo de
$h^{-1}(Y')$, y por
$\widehat{g}:\widehat{Z}\to\widehat{X}$ la extensión de $g$ a
las compleciones.

El módulo $\widehat{g}^{\,*}(\sh{F})$ es un
$\sh{O}_{\widehat{Z}}$-módulo coherente cuyo soporte está contenido en
$g^{-1}(E)$ y es, por tanto, propio sobre $Y$, pues $g$ es
proyectivo y, por tanto, propio \sref[II]{II.5.4.6}. Como $h$ es
cuasiproyectivo, $\widehat{g}^{\,*}(\sh{F})$ es algebraizable por
\sref{III.5.2.6}.
\oldpage[III]{156}
Se deduce de \sref{III.5.3.3}, puesto que $g$ es propio, que
\[
  \widehat{g}_*
  \bigl(\widehat{g}^{\,*}(\sh{F})\bigr)
\]
es un $\sh{O}_{\widehat{X}}$-módulo algebraizable.

Apliquemos \sref{III.5.3.4} a $\sh{F}$ y $g$. El núcleo $\sh{P}$
y el conúcleo $\sh{R}$ de
\[
  \rho_{\sh{F}}:
  \sh{F}\longrightarrow
  \widehat{g}_*\bigl(\widehat{g}^{\,*}(\sh{F})\bigr)
\]
son anulados por una potencia $\widehat{\sh{M}}^{\,n}$. Sea $T$
el subpreesquema cerrado de $X$ definido por $\sh{M}^n$, cuyo
espacio subyacente es $X-U$, y sea $j:T\to X$ la inmersión
canónica. Su extensión
$\widehat{j}:\widehat{T}\to\widehat{X}$ es la inmersión canónica
\sref[I]{I.10.14.7}, y podemos escribir
\[
  \sh{P}=
  \widehat{j}_*\bigl(\widehat{j}^{\,*}(\sh{P})\bigr),
  \qquad
  \sh{R}=
  \widehat{j}_*\bigl(\widehat{j}^{\,*}(\sh{R})\bigr).
\]
Como $U$ no es vacío, la hipótesis de inducción muestra que
$\widehat{j}^{\,*}(\sh{P})$ y
$\widehat{j}^{\,*}(\sh{R})$ son $\sh{O}_{\widehat{T}}$-módulos
algebraizables. El Lema \sref{III.5.3.3} muestra entonces
que $\sh{P}$ y $\sh{R}$ son algebraizables, y
\sref{III.5.3.2} demuestra finalmente que $\sh{F}$ es algebraizable.
\end{env}


\subsection{Aplicación: comparación entre morfismos de preesquemas ordinarios y morfismos de preesquemas formales. Preesquemas formales algebraizables}
\label{subsection:III.5.4}

\begin{theorem}[5.4.1]
\label{III.5.4.1}
Sea $A$ un anillo noetheriano ádico, sea $\mathfrak{J}$ un ideal
de definición de $A$, y pongamos
\[
  S=\Spec(A),\qquad S'=\mathrm{V}(\mathfrak{J}).
\]
Sea $u:X\to S$ un morfismo propio y sea $v:Y\to S$ un
morfismo separado de tipo finito. Sean
$\widehat{S}$, $\widehat{X}$ y $\widehat{Y}$ las compleciones
de $S$, $X$ e $Y$ a lo largo de $S'$, $u^{-1}(S')$ y $v^{-1}(S')$,
respectivamente. Para todo $S$-morfismo $f:X\to Y$, sea
$\widehat{f}:\widehat{X}\to\widehat{Y}$ su extensión a las
compleciones. Entonces la aplicación $f\mapsto\widehat{f}$ es una biyección
\[
  \Hom_S(X,Y)
  \xrightarrow{\ \sim\ }
  \Hom_{\widehat{S}}(\widehat{X},\widehat{Y}).
\]
\end{theorem}

\begin{proof}
Demostremos primero la inyectividad. Supongamos que $f,g:X\to Y$ son
$S$-morfismos tales que $\widehat{f}=\widehat{g}$. Existe entonces
un entorno abierto $V$ de $X'=u^{-1}(S')$ sobre el cual $f$ y $g$
coinciden \sref[I]{I.10.9.4}. Como $u$ es cerrado,
\sref{III.5.1.3.1} da $V=X$, y por tanto $f=g$.

Demostremos ahora la sobreyectividad. Sea
$h:\widehat{X}\to\widehat{Y}$ un
$\widehat{S}$-morfismo. Pongamos
\[
  Z=X\times_S Y,
\]
y denotemos las proyecciones canónicas por $p:Z\to X$ y
$q:Z\to Y$. El preesquema $Z$ es de tipo finito sobre $S$
\sref[I]{I.6.3.4} y, por tanto, noetheriano. Sea
$\widehat{Z}$ su compleción a lo largo de
\[
  Z'=p^{-1}\bigl(u^{-1}(S')\bigr).
\]
El preesquema formal $\widehat{Z}$ se identifica canónicamente con
$\widehat{X}\times_{\widehat{S}}\widehat{Y}$, y sus proyecciones
sobre $\widehat{X}$ y $\widehat{Y}$ se identifican con
$\widehat{p}$ y $\widehat{q}$ \sref[I]{I.10.9.7}.

Como $Y$ es separado sobre $S$, $\widehat{Y}$ es separado sobre
$\widehat{S}$ \sref[I]{I.10.15.7}. El morfismo grafo
\[
  \Gamma_h=(1_{\widehat{X}},h):
  \widehat{X}\longrightarrow\widehat{Z}
\]
es, por consiguiente, una inmersión cerrada \sref[I]{I.10.15.4}. Sea
$\mathfrak{T}$ el subpreesquema formal cerrado de
$\widehat{Z}$ asociado a esta inmersión, y sea
$j:\mathfrak{T}\to\widehat{Z}$ la inmersión canónica. Así,
\[
  \Gamma_h=j\circ w,
\]
donde $w:\widehat{X}\to\mathfrak{T}$ es un isomorfismo
\sref[I]{I.10.14.3} cuyo inverso es
$\widehat{p}\circ j$. Además, $\mathfrak{T}$ es propio sobre
$\widehat{S}$ porque $\widehat{X}$ lo es. Se deduce de
\sref{III.5.1.8} que existe un subpreesquema cerrado $T$ de $Z$
tal que
\[
  \mathfrak{T}
  =\widehat{T}
  =T_{/(T\cap Z')},
\]
y que $j=\widehat{i}$, donde $i:T\to Z$ es la inmersión
canónica \sref[I]{I.10.14.7}.

El morfismo $p\circ i:T\to X$ es un isomorfismo: su extensión a
las compleciones es el isomorfismo
$\widehat{p}\circ\widehat{i}$, y basta aplicar
\sref{III.4.6.8}, observando como antes que $S$ es el único
entorno de $S'$ en $S$.
\oldpage[III]{157}
Sea $g:X\to T$ el inverso de $p\circ i$, y pongamos
\[
  f=q\circ i\circ g.
\]
Entonces $f:X\to Y$ tiene por grafo $\Gamma_f=i\circ g$. Como
$\widehat{g}$ es el inverso de
$\widehat{p\circ i}$, se tiene
\[
  \widehat{\Gamma_f}
  =\widehat{i}\circ\widehat{g}
  =j\circ w
  =\Gamma_h.
\]
Pero $\widehat{\Gamma_f}=\Gamma_{\widehat{f}}$
\sref[I]{I.10.9.8}; por tanto, $h=\widehat{f}$, lo que termina la
demostración.
\end{proof}

En lenguaje categórico, el funtor
$X\mapsto\widehat{X}$ es, por tanto, plenamente fiel
\sref[0]{0.8.1.6} de la categoría de los preesquemas propios sobre
$\Spec(A)$ en la categoría de los preesquemas formales propios sobre
$\Spf(A)$, para todo anillo noetheriano ádico $A$. Proporciona, por
consiguiente, una equivalencia entre la primera categoría y una
subcategoría plena de la segunda. Los objetos de esta subcategoría se
denominarán en lo sucesivo \emph{preesquemas formales algebraizables}.
Para tal preesquema formal $\mathfrak{X}$, existe un preesquema
ordinario $X$, propio sobre $\Spec(A)$ y determinado salvo isomorfismo
único, tal que $\mathfrak{X}\simeq\widehat{X}$.

\begin{env}[5.4.2]
\label{III.5.4.2}
\emph{Escolio.}
Con la notación de \sref{III.5.4.1}, pongamos
\[
  S_n=\Spec(A/\mathfrak{J}^{n+1}),\qquad
  X_n=X\times_S S_n,\qquad
  Y_n=Y\times_S S_n.
\]
Se deduce de \sref{III.5.4.1} y \sref[I]{I.10.12.3} que
dar un $S$-morfismo $f:X\to Y$ equivale a dar un sistema inductivo
ádico $(S_n)$ \sref[I]{I.10.12.2} de
$S_n$-morfismos $f_n:X_n\to Y_n$.
\end{env}

\begin{remark}[5.4.3]
\label{III.5.4.3}
Contrariamente a lo que podría sugerir el teorema de existencia
\sref{III.5.1.6}, existen preesquemas formales propios sobre $\Spf(A)$
que \emph{no son algebraizables}, así como existen espacios analíticos
compactos que no proceden de variedades algebraicas complejas. Más
adelante encontraremos tales preesquemas en la «teoría de moduli»,
que, cuando el cuerpo base es $\mathbb{C}$, trata precisamente de las
variaciones infinitesimales de la estructura compleja de una variedad
algebraica completa; es bien sabido que tales variaciones pueden dar
lugar a variedades analíticas que no son algebraicas.
\end{remark}

\begin{proposition}[5.4.4]
\label{III.5.4.4}
Sea $A$ un anillo noetheriano ádico, pongamos
$\mathfrak{S}=\Spf(A)$, y sean
\[
  g:\mathfrak{X}\to\mathfrak{S},
  \qquad
  h:\mathfrak{Y}\to\mathfrak{S}
\]
dos morfismos propios de preesquemas formales. Sea
$f:\mathfrak{X}\to\mathfrak{Y}$ un
$\mathfrak{S}$-morfismo. Si $f$ es finito y
$\mathfrak{Y}$ es algebraizable, entonces $\mathfrak{X}$ es
algebraizable.
\end{proposition}

\begin{proof}
Las hipótesis sobre $g$ y $h$ implican ya que $f$ es propio
\sref{III.3.4.1}; para que $f$ sea finito, basta que la fibra
$f^{-1}(y)$ sea finita para todo $y\in\mathfrak{Y}$
\sref[0]{0.4.8.11}. Por hipótesis,
\[
  \sh{B}=f_*(\sh{O}_{\mathfrak{X}})
\]
es una $\sh{O}_{\mathfrak{Y}}$-álgebra coherente
\sref[0]{0.4.8.6}. Escribamos
$\mathfrak{Y}=\widehat{Y}$ y $h=\widehat{w}$, donde
$w:Y\to\Spec(A)$ es un morfismo propio de preesquemas ordinarios. El
teorema de existencia proporciona una $\sh{O}_Y$-álgebra coherente $\sh{C}$
tal que $\sh{B}=\widehat{\sh{C}}$. Pongamos
\[
  X=\Spec(\sh{C}),
\]
y sea $u:X\to Y$ el morfismo estructural. La definición de
$\mathfrak{X}$ a partir de $\sh{B}$ \sref[0]{0.4.8.7} muestra entonces que
$\mathfrak{X}$ es canónicamente isomorfo a $\widehat{X}$ y
que $f$ se identifica con $\widehat{u}$; basta comprobarlo
cuando $Y$ es afín.
\end{proof}

La Proposición \sref{III.5.1.8} es un caso particular de
\sref{III.5.4.4}.

\begin{theorem}[5.4.5]
\label{III.5.4.5}
Sea $A$ un anillo noetheriano ádico, sea $\mathfrak{J}$ un ideal
de definición de $A$, pongamos
\[
  S=\Spec(A),\qquad
  \mathfrak{S}=\widehat{S}=\Spf(A),
\]
y sea $f:\mathfrak{X}\to\mathfrak{S}$ un morfismo propio de
preesquemas formales. Pongamos
\[
  S_k=\Spec(A/\mathfrak{J}^{k+1}),\qquad
  X_k=\mathfrak{X}\times_{\mathfrak{S}}S_k,
\]
y, para todo $\sh{O}_{\mathfrak{X}}$-módulo $\sh{F}$, pongamos
\[
  \sh{F}_k
  =\sh{F}\otimes_{\sh{O}_{\mathfrak{X}}}\sh{O}_{X_k}
  =\sh{F}/\mathfrak{J}^{k+1}\sh{F}.
\]
\oldpage[III]{158}
Sea $\sh{L}$ un
$\sh{O}_{\mathfrak{X}}$-módulo invertible, y supongamos que
\[
  \sh{L}_0=\sh{L}/\mathfrak{J}\sh{L}
\]
es un $\sh{O}_{X_0}$-módulo amplio. Entonces $\mathfrak{X}$ es
algebraizable. Además, si $X$ es un $S$-preesquema propio tal
que $\mathfrak{X}\simeq\widehat{X}$, existe un
$\sh{O}_X$-módulo amplio $\sh{M}$ tal que
$\sh{L}\simeq\widehat{\sh{M}}$; en particular, $X$ es proyectivo
sobre $S$.
\end{theorem}

\begin{proof}
Apliquemos \sref{III.5.2.3} con
$\sh{F}=\sh{O}_{\mathfrak{X}}$. Existe un entero $n_0$ tal
que, para $n\geq n_0$, el homomorfismo canónico
\[
  \Gamma(\mathfrak{X},\sh{L}^{\otimes n})
  \longrightarrow
  \Gamma(X_0,\sh{L}_0^{\otimes n})
\]
es sobreyectivo. Elijamos $n\geq n_0$ suficientemente grande para que
$\sh{L}_0^{\otimes n}$ sea muy amplio para $S_0$
\sref[II]{II.4.5.10}. Como
$f_0:X_0\to S_0$ es propio,
$\Gamma(X_0,\sh{L}_0^{\otimes n})$ es un $A$-módulo de tipo finito
\sref{III.3.2.1}. Existe, por consiguiente, un
$A$-submódulo de tipo finito
\[
  E\subset\Gamma(\mathfrak{X},\sh{L}^{\otimes n})
\]
cuya imagen en $\Gamma(X_0,\sh{L}_0^{\otimes n})$ es el módulo
entero.

Para todo $k\geq0$, consideremos el homomorfismo
\[
  u_k:
  E/\mathfrak{J}^{k+1}E
  \longrightarrow
  \Gamma(X_k,\sh{L}_k^{\otimes n})
\]
inducido por la inyección canónica
$E\to\Gamma(\mathfrak{X},\sh{L}^{\otimes n})$. El módulo
$(f_k)_*(\sh{L}_k^{\otimes n})$ es cuasi-coherente, y
\[
  \Gamma\bigl(S_k,(f_k)_*(\sh{L}_k^{\otimes n})\bigr)
  =
  \Gamma(X_k,\sh{L}_k^{\otimes n}).
\]
Así, $u_k$ define un homomorfismo
\[
  \widetilde{u}_k:
  \widetilde{E/\mathfrak{J}^{k+1}E}
  \longrightarrow
  (f_k)_*(\sh{L}_k^{\otimes n}),
\]
y, por tanto, por adjunción, un homomorfismo
\[
  \widetilde{u}_k^{\,\sharp}:
  f_k^*\bigl(\widetilde{E/\mathfrak{J}^{k+1}E}\bigr)
  \longrightarrow
  \sh{L}_k^{\otimes n}.
\]
Pongamos
\[
  \sh{G}_k=
  f_k^*\bigl(\widetilde{E/\mathfrak{J}^{k+1}E}\bigr).
\]
Se tiene
$\sh{G}_0=\sh{G}_k/\mathfrak{J}\sh{G}_k$
\sref[I]{I.9.1.5}, y
\[
  \widetilde{u}_0^{\,\sharp}:
  \sh{G}_0/\mathfrak{J}\sh{G}_0
  \longrightarrow
  \sh{L}_0^{\otimes n}/
  \mathfrak{J}\sh{L}_0^{\otimes n}
\]
se obtiene de $\widetilde{u}_k^{\,\sharp}$ pasando a los
cocientes.

Por la definición de $E$, $\widetilde{u}_0^{\,\sharp}$ es el
homomorfismo canónico
\[
  \sigma:
  f_0^*\bigl((f_0)_*(\sh{L}_0^{\otimes n})\bigr)
  \longrightarrow
  \sh{L}_0^{\otimes n}.
\]
Como $\sh{L}_0^{\otimes n}$ es muy amplio,
$\widetilde{u}_0^{\,\sharp}$ es sobreyectivo
\sref[II]{II.4.4.3}. El lema de Nakayama muestra entonces que todo
$\widetilde{u}_k^{\,\sharp}$ es sobreyectivo.

Cada $\widetilde{u}_k^{\,\sharp}$ define, por tanto,
\sref[II]{II.4.2.2} un $S_k$-morfismo
\[
  g_k:X_k\longrightarrow
  P_k=\mathbf{P}(E/\mathfrak{J}^{k+1}E).
\]
Para $h\leq k$,
\[
  P_h=P_k\times_{S_k}S_h
\]
por \sref[II]{II.4.1.3}. Las relaciones entre los
$\widetilde{u}_k^{\,\sharp}$ y \sref[II]{II.4.2.10} muestran que
$(g_k)$ es un sistema inductivo ádico $(S_k)$
\sref[I]{I.10.12.2}. Los $g_k$ definen, por consiguiente, un
$\mathfrak{S}$-morfismo de preesquemas formales
\[
  g:\mathfrak{X}\longrightarrow\mathfrak{P},
\]
donde $\mathfrak{P}$ es el límite inductivo del sistema $(P_k)$,
o, equivalentemente, la compleción $\widehat{P}$ de
$P=\mathbf{P}(E)$. El hecho de que
$\sh{L}_0^{\otimes n}$ sea muy amplio implica que $g_0$ es una inmersión cerrada
\sref[II]{II.4.4.3}. Por tanto, $g$ es una inmersión cerrada de
preesquemas formales \sref[0]{0.4.8.10}, y $\mathfrak{X}$ es algebraizable
por \sref{III.5.1.8}.

El teorema de existencia \sref{III.5.1.6} muestra ahora que $\sh{L}$ es
la compleción $\widehat{\sh{M}}$ de un
$\sh{O}_X$-módulo invertible $\sh{M}$. Además,
$\sh{L}^{\otimes n}$ es la compleción de
$\sh{M}^{\otimes n}$ \sref[I]{I.10.8.10}, y los homomorfismos
$\widetilde{u}_k^{\,\sharp}$ definen un único homomorfismo
\[
  v:f^*(\widetilde{E})
  \longrightarrow
  \sh{M}^{\otimes n}
\]
\sref{III.5.1.7}. Como los
$\widetilde{u}_k^{\,\sharp}$ son sobreyectivos,
$\widehat{v}$ es sobreyectivo \sref[I]{I.10.11.5}, y por tanto también
lo es $v$ \sref{III.5.1.3}. Sea $r:X\to P$ el morfismo definido
por $v$ \sref[II]{II.4.2.2}. Su extensión a las compleciones es
$g$. Como $g$ es una inmersión cerrada, también lo es $r$, por
\sref{III.5.1.8} y \sref{III.5.4.1}. Se deduce que
$\sh{M}^{\otimes n}$ es muy amplio \sref[II]{II.4.4.6} y, por tanto,
que $\sh{M}$ es amplio \sref[II]{II.4.5.10}.
\end{proof}

\begin{remark}[5.4.6]
\label{III.5.4.6}
Sea $A$ un anillo noetheriano ádico, pongamos
$\mathfrak{S}=\Spf(A)$, y sea
$f:\mathfrak{X}\to\mathfrak{S}$ un morfismo propio de preesquemas
formales. Sea $\sh{N}$ el ideal coherente de
$\sh{O}_{\mathfrak{X}}$ tal que, para todo subconjunto abierto
formal afín $U$ de $\mathfrak{X}$,
\[
  \Gamma(U,\sh{N})
\]
es el nilradical de
$\Gamma(U,\sh{O}_{\mathfrak{X}})$. La existencia de este ideal
se deduce fácilmente de \sref[I]{I.10.10.2} y del hecho de que
todo homomorfismo de anillos $B\to C$ envía el nilradical de $B$ en
el de $C$. Sea $\mathfrak{X}'$ el subpreesquema formal cerrado de
$\mathfrak{X}$ definido por $\sh{N}$
\sref[I]{I.10.14.2}. Sería interesante saber si la
algebraizabilidad de $\mathfrak{X}'$ implica la algebraizabilidad del
propio $\mathfrak{X}$.

Sin duda se llegaría a una solución de este problema si se supiera
clasificar, por ejemplo mediante invariantes
\oldpage[III]{159}
de naturaleza cohomológica, las \emph{extensiones} de un haz estructural
$\sh{O}_{\mathfrak{X}}$---para un preesquema ordinario o
formal---por un ideal de cuadrado nulo. Equivalentemente, estas son las
$\sh{O}_{\mathfrak{X}}$-álgebras $\sh{A}$ tales que
$\sh{O}_{\mathfrak{X}}\simeq\sh{A}/\sh{J}$, donde $\sh{J}$ es un
ideal de cuadrado nulo de $\sh{A}$.
\end{remark}


\subsection{Una descomposición de ciertos preesquemas}
\label{subsection:III.5.5}

\begin{proposition}[5.5.1]
\label{III.5.5.1}
Sea $A$ un anillo noetheriano ádico, sea $\mathfrak{J}$ un ideal
de definición de $A$, y pongamos $Y=\Spec(A)$. Sea $f:X\to Y$ un
morfismo separado de tipo finito, y pongamos
\[
  Y_0=\Spec(A/\mathfrak{J}),\qquad
  X_0=X\times_Y Y_0=f^{-1}(Y_0).
\]
Sea $Z_0$ un subconjunto abierto de $X_0$ que es propio sobre $Y_0$.
Entonces existe un subconjunto abierto y cerrado $Z$ de $X$, propio sobre
$Y$, tal que
\[
  Z\cap X_0=Z_0.
\]
\end{proposition}

\begin{proof}
Por hipótesis, existe un subconjunto abierto $T$ de $X$ tal que
$T\cap X_0=Z_0$. Sea $\widehat{T}$ la compleción a lo largo de $Z_0$
del preesquema inducido por $X$ sobre $T$. El soporte de
$\sh{O}_{\widehat{T}}$ es $Z_0$, que es propio sobre $Y_0$;
por tanto, $\widehat{T}$ es propio sobre
$\widehat{Y}=\Spf(A)$ \sref{III.3.4.1}.

Se deduce de \sref{III.5.1.8} que existe un subpreesquema cerrado
$Z$ de $T$, propio sobre $Y$, tal que, si
$i:Z\to T$ es la inmersión canónica, entonces
\[
  \widehat{i}:\widehat{Z}\longrightarrow\widehat{T}
\]
es un isomorfismo; aquí $\widehat{Z}$ es la compleción de $Z$
a lo largo de $Z_0$. Por \sref{III.4.6.8}, existe un entorno abierto
$V$ de $Z_0$ en $T$ tal que
$i^{-1}(V)\to V$ es un isomorfismo. Pero $i^{-1}(V)$ es un
entorno de $Z_0$ en $Z$ y, por tanto, es igual a $Z$ por
\sref{III.5.1.3.1}. Así, $Z$ es abierto en $T$, y por tanto en $X$.
También es cerrado en $X$, pues $Z$ es propio sobre $Y$ y $X$ es
separado sobre $Y$. Esto demuestra la proposición.
\end{proof}

\begin{corollary}[5.5.2]
\label{III.5.5.2}
Si $X_0$ es propio sobre $Y_0$, entonces $X$ es la unión de dos
subconjuntos abiertos disjuntos $Z$ y $Z'$ tales que $Z$ es propio sobre
$Y$ y contiene a $X_0$. Además, todo subconjunto cerrado $P$ de $X$
que es propio sobre $Y$ está contenido en $Z$.
\end{corollary}

\begin{proof}
Apliquemos \sref{III.5.5.1} con $Z_0=X_0$, y pongamos $Z'=X-Z$. Para la
última afirmación, $P\cap Z'$ es cerrado en $P$ y, por tanto, propio
sobre $Y$. Si no fuera vacío, $f(P\cap Z')$ sería un subconjunto cerrado
no vacío de $Y$ y, por tanto, cortaría a $Y_0$
\sref{III.5.1.3.1}. Esto contradice
$Z'\cap X_0=\varnothing$ y demuestra que $P\subset Z$.
\end{proof}

\bigskip
\noindent\emph{(Continuará.)}

% La página impresa 504 está en blanco.
\clearpage
\thispagestyle{empty}
\null
\clearpage
\thispagestyle{plain}

\oldpage[III]{161}
\section*{BIBLIOGRAFÍA (continuación)}

\begin{flushleft}
[27] P. Gabriel, \emph{Des catégories abéliennes}, tesis, París (1961).

\medskip
[28] J. E. Roos, \emph{Sur les foncteurs dérivés de $\varprojlim$.
Applications}, C. R. Acad. Sci. Paris, vol. CCLII, 1961, pp. 3702--3704.

\medskip
[29] H. Cartan, \emph{Séminaire de l'École Normale Supérieure},
13.er año (1960--1961), exposición n.º 11.
\end{flushleft}

\clearpage
\oldpage[III]{162}
\section*{ÍNDICE DE NOTACIONES}

\begin{flushleft}
$\Set$: 0, 8.1.1.

$h_X(Y)$, $h_X(u)$ ($X$, $Y$ objetos de una categoría $\C$, $u$ un morfismo
de $\C$): 0, 8.1.1.

$h_w(Y)$ ($Y$ un objeto y $w$ un morfismo de $\C$): 0, 8.1.2.

$Z_k(E_r^{pq})$, $B_k(E_r^{pq})$ ($k\geq r$): 0, 11.1.1.

$Z_\infty(E_2^{pq})$, $B_\infty(E_2^{pq})$: 0, 11.1.1.

$K^\bullet$, $K_\bullet$ (complejos): 0, 11.2.1.

$K^{\bullet\bullet}$, $K_{\bullet\bullet}$ (bicomplejos): 0, 11.2.1.

$B^i(K^\bullet)$, $Z^i(K^\bullet)$, $H^i(K^\bullet)$,
$H^\bullet(K^\bullet)$: 0, 11.2.1.

$B_i(K_\bullet)$, $Z_i(K_\bullet)$, $H_i(K_\bullet)$,
$H_\bullet(K_\bullet)$: 0, 11.2.1.

$T(K^\bullet)$, $T(K_\bullet)$ ($T$ un funtor): 0, 11.2.1.

$E(K^\bullet)$ ($K^\bullet$ un complejo filtrado): 0, 11.2.2.

$K^{i,\bullet}$, $K^{\bullet,j}$, $Z_{\mathrm{II}}^p(K^{i,\bullet})$,
$B_{\mathrm{II}}^p(K^{i,\bullet})$, $H_{\mathrm{II}}^p(K^{i,\bullet})$,
$Z_{\mathrm I}^p(K^{\bullet,j})$, $B_{\mathrm I}^p(K^{\bullet,j})$,
$H_{\mathrm I}^p(K^{\bullet,j})$: 0, 11.3.1.

$H_{\mathrm{II}}^p(K^{\bullet\bullet})$,
$Z_{\mathrm I}^q(H_{\mathrm{II}}^p(K^{\bullet\bullet}))$,
$B_{\mathrm I}^q(H_{\mathrm{II}}^p(K^{\bullet\bullet}))$,
$H_{\mathrm I}^q(H_{\mathrm{II}}^p(K^{\bullet\bullet}))$,
$H^n(K^{\bullet\bullet})$: 0, 11.3.1.

$F_{\mathrm I}^p(K^{\bullet\bullet})$,
$F_{\mathrm{II}}^p(K^{\bullet\bullet})$, ${}'E(K^{\bullet\bullet})$,
${}''E(K^{\bullet\bullet})$: 0, 11.3.2.

$K_{i,\bullet}$, $K_{\bullet,j}$, $H_p^{\mathrm{II}}(K_{i,\bullet})$,
$H_p^{\mathrm I}(K_{\bullet,j})$, $H_p^{\mathrm{II}}(K_{\bullet\bullet})$,
$H_p^{\mathrm I}(K_{\bullet\bullet})$,
$H_q^{\mathrm I}(H_p^{\mathrm{II}}(K_{\bullet\bullet}))$,
$H_q^{\mathrm{II}}(H_p^{\mathrm I}(K_{\bullet\bullet}))$,
$H_n(K_{\bullet\bullet})$: 0, 11.3.5.

$\RR^\bullet T(K^\bullet)$ ($T$ un funtor covariante aditivo): 0, 11.4.3.

$\RR^\bullet T(K^\bullet,K'^{\bullet})$ ($T$ un bifuntor covariante
aditivo): 0, 11.4.6.

$\LL_\bullet T(K_\bullet)$ ($T$ un funtor covariante aditivo): 0, 11.6.2.

$\LL_\bullet T(K_\bullet,K'_\bullet)$ ($T$ un bifuntor covariante
aditivo): 0, 11.6.6.

$\LL_\bullet T(K_{\bullet\bullet})$ ($T$ un funtor covariante aditivo):
0, 11.7.2.

$|\sigma|$, $\Sigma(A)$, $\mathrm C_\bullet(A)$, $\mathrm L_\bullet(A)$
($A$ un conjunto finito, $\sigma$ un símplex de $A$): 0, 11.8.1.

$\mathrm C^\bullet(A,B;\sh{S})$, $\mathrm L^\bullet(A,B;\sh{S})$
($A$, $B$ conjuntos finitos, $\sh{S}$ un sistema de coeficientes): 0, 11.8.4.

$\mathrm P^\bullet(A,B;\sh{S})$ ($A$, $B$ conjuntos finitos, $\sh{S}$ un
sistema de coeficientes): 0, 11.8.6.

$\mathrm C^\bullet(A,B;\sh{S}^\bullet)$,
$\mathrm L^\bullet(A,B;\sh{S}^\bullet)$ ($A$, $B$ conjuntos finitos,
$\sh{S}^\bullet$ un complejo de sistemas de coeficientes): 0, 11.8.8.

$\mathrm P^\bullet(A,B;\sh{S}^\bullet)$ ($A$, $B$ conjuntos finitos,
$\sh{S}^\bullet$ un complejo de sistemas de coeficientes): 0, 11.8.10.

$\chi(M^\bullet)$ ($M^\bullet$ un complejo finito de módulos de longitud finita):
0, 11.10.1.

$\check C^\bullet(\mathfrak U,\sh{F})$ ($\sh{F}$ un haz de grupos abelianos
sobre $X$, $\mathfrak U$ un recubrimiento abierto de $X$): 0, 12.1.2.

$\RR^p f_*(\sh{F})$ ($f:X\to Y$ un morfismo de espacios anillados,
$\sh{F}$ un haz de $\sh{O}_X$-módulos): 0, 12.2.1.

$\sh{H}^p(f,\sh{K}^\bullet)$,
$\sh{H}^p_f(\sh{K}^\bullet)$, ${}'E(f,\sh{K}^\bullet)$,
${}''E(f,\sh{K}^\bullet)$ ($f:X\to Y$ un morfismo de espacios anillados,
$\sh{K}^\bullet$ un complejo de $\sh{O}_X$-módulos): 0, 12.4.1.

$\mathbf H^\bullet(X,\sh{K}^\bullet)$,
$\mathbf H^\bullet(V,\sh{K}^\bullet)$ ($X$ un espacio anillado,
$\sh{K}^\bullet$ un complejo de $\sh{O}_X$-módulos, $V$ abierto en $X$):
0, 12.4.2.

$\operatorname{gr}^\bullet(A)$ ($A$ un sistema proyectivo que satisface (ML)):
0, 13.4.2.

$E(A)$ ($A$ un objeto filtrado): 0, 13.6.4.

$K_\bullet(\mathbf f)$, $i_{\mathbf f}$ ($\mathbf f$ una familia finita de
elementos de un anillo): III, 1.1.1.

$K^\bullet(\mathbf f,M)$, $K_\bullet(\mathbf f,M)$ ($\mathbf f$ una familia
finita de elementos de un anillo $A$, $M$ un $A$-módulo): III, 1.1.2.

$H^\bullet(\mathbf f,M)$, $H_\bullet(\mathbf f,M)$ ($\mathbf f$ una familia
finita de elementos de un anillo $A$, $M$ un $A$-módulo): III, 1.1.3.

$C^\bullet((\mathbf f),M)$, $H^\bullet((\mathbf f),M)$ ($\mathbf f$ una
familia finita de elementos de un anillo $A$, $M$ un $A$-módulo): III, 1.1.6.

$H^\bullet(U,\sh{F}^{(*)})$,
$H^\bullet(\mathfrak U,\sh{F}^{(*)})$ ($\sh{F}$ un
$\sh{O}_X$-módulo cuasi-coherente, $U$ abierto en $X$, $\mathfrak U$ un recubrimiento abierto de $X$):
III, 2.1.1.
\end{flushleft}

\clearpage
\oldpage[III]{163}
\section*{ÍNDICE TERMINOLÓGICO}

\begin{flushleft}
Término de convergencia de una sucesión espectral: 0, 11.1.1.

Algebraizable ($\sh{O}_X$-módulo): III, 5.2.1.

Algebraizable (esquema formal): III, 5.4.2.

Analíticamente íntegro (anillo): III, 4.3.6.

Aplicación cuasicompacta: 0, 9.1.1.

Aumentación de una resolución: 0, 11.4.1.

Característica de Euler--Poincaré de un complejo: 0, 11.10.1.

Característica de Euler--Poincaré de un $\sh{O}_X$-módulo coherente
($X$ proyectivo sobre $\Spec(A)$, $A$ un anillo artiniano): III, 2.5.1.

Cocadena bialternada: 0, 11.8.4.

Complejo definido por un bicomplejo: 0, 11.3.1.

Complejo de álgebra exterior: III, 1.1.1.

Condición (ML): 0, 13.1.2.

Constructible (subconjunto, conjunto): 0, 9.1.2.

Constructible (función): 0, 9.3.1.

Producto cup: 0, 12.1.2.

Dihomomorfismo para estructuras algebraicas sobre categorías: 0, 8.2.1.

Exacto (subconjunto) en una categoría abeliana: III, 3.1.1.

Esencialmente constante (sistema proyectivo): 0, 13.4.2.

Factorización de Stein: III, 4.3.3.

Codiscreta, coseparada, discreta, exhaustiva, finita, separada
(filtración): 0, 11.1.3.

Final (objeto) de una categoría: 0, 8.1.10.

Finito (morfismo) de preesquemas formales: III, 4.8.2.

Finito (preesquema formal) sobre $\mathfrak Y$,
$\mathfrak Y$-finito (preesquema formal): III, 4.8.2.

Funtor covariante canónico $\C\to\CHom(\C\op,\Set)$: 0, 8.1.2.

Género aritmético: III, 2.5.1.

Geométricamente conexa (fibra): III, 4.3.4.

Homomorfismo de estructuras algebraicas sobre categorías: 0, 8.2.1.

Hipercohomología de un funtor respecto de un complejo $K^\bullet$:
0, 11.4.3.

Hipercohomología de un bifuntor respecto de dos complejos
$K^\bullet$, $K'^{\bullet}$: 0, 11.4.6.

Hiperhomología de un funtor respecto de un complejo $K_\bullet$:
0, 11.6.2.

Hiperhomología de un bifuntor respecto de dos complejos $K_\bullet$,
$K'_\bullet$: 0, 11.6.6.

Hiperhomología de un funtor respecto de un bicomplejo: 0, 11.7.2.

Localmente constructible (conjunto): 0, 9.1.11.

Ley de composición externa (interna) sobre los objetos de una categoría: 0, 8.2.1.

\oldpage[III]{164}

$S$-$\C$-módulo filtrado, $\operatorname{gr}^\bullet(S)$-$\C$-módulo graduado,
$\operatorname{gr}^{\bullet\bullet}(S)$-$\C$-módulo bigraduado: 0, 13.6.6.

Morfismo de sucesiones espectrales: 0, 11.1.2.

Número geométrico de componentes conexas de una fibra: III, 4.3.4.

Objeto grupo en $\C$, $\C$-grupo, $\C$-anillo, $\C$-módulo: 0, 8.2.3.

Objeto de bordes, objeto de cobordes, objeto de cociclos, objeto de
ciclos: 0, 11.2.1.

Objeto de imágenes universales: 0, 13.1.1.

Objeto graduado asociado a un sistema proyectivo que satisface (ML): 0, 13.4.2.

Objeto graduado acotado inferiormente (superiormente): 0, 11.2.1.

Subconjunto propio (sobre $\mathfrak Y$) de un preesquema formal: III, 3.4.1.

Plena (subcategoría): 0, 8.1.5.

Plenamente fiel (funtor): 0, 8.1.5.

Polinomio de Hilbert: III, 2.5.3.

Propio (morfismo) de preesquemas formales: III, 3.4.1.

Representable (funtor): 0, 8.1.8.

Cohomológica, derecha, izquierda, homológica, inyectiva, libre, plana, proyectiva
(resolución): 0, 11.4.1.

Resolución de Cartan--Eilenberg derecha, inyectiva: 0, 11.4.2.

Resolución de Cartan--Eilenberg izquierda, proyectiva: 0, 11.6.1.

Retrocompacto (conjunto): 0, 9.1.1.

Estricto (sistema proyectivo): 0, 13.4.2.

Sucesión espectral en una categoría abeliana: 0, 11.1.1.

Débilmente convergente, regular, coregular, birregular (sucesión espectral):
0, 11.1.3.

Sucesión espectral degenerada: 0, 11.1.6.

Sucesión espectral de un funtor relativa a un objeto filtrado: 0, 13.6.4.

Sucesiones espectrales de un bicomplejo: 0, 11.3.2 y 11.3.5.

Sucesiones espectrales de hipercohomología de un funtor respecto de un complejo:
0, 11.4.3.

Sucesiones espectrales de hiperhomología de un funtor respecto de un complejo:
0, 11.6.2.

Sucesiones espectrales de hiperhomología de un funtor respecto de un bicomplejo:
0, 11.7.2.

Sistema de coeficientes: 0, 11.8.4.

Unirrama (anillo, punto): III, 4.3.6.

Universalmente abierto (morfismo): III, 4.3.9.
\end{flushleft}

\clearpage
\oldpage[III]{165}
\section*{ÍNDICE GENERAL}

\begin{flushleft}
\textsc{Capítulo 0.} --- \textbf{Preliminares (continuación)} \dotfill 5

\medskip
\S~8. Funtores representables \dotfill 5

8.1. Funtores representables \dotfill 5

8.2. Estructuras algebraicas en categorías \dotfill 9

\medskip
\S~9. Conjuntos constructibles \dotfill 12

9.1. Conjuntos constructibles \dotfill 12

9.2. Subconjuntos constructibles de espacios noetherianos \dotfill 14

9.3. Funciones constructibles \dotfill 16

\medskip
\S~10. Suplemento sobre los módulos planos \dotfill 17

10.1. Relaciones entre módulos planos y módulos libres \dotfill 17

10.2. Criterios locales de planitud \dotfill 18

10.3. Existencia de extensiones planas de anillos locales \dotfill 20

\medskip
\S~11. Suplemento sobre álgebra homológica \dotfill 23

11.1. Recordatorio sobre las sucesiones espectrales \dotfill 23

11.2. La sucesión espectral de un complejo filtrado \dotfill 27

11.3. Las sucesiones espectrales de un bicomplejo \dotfill 29

11.4. Hipercohomología de un funtor respecto de un complejo $K^\bullet$
\dotfill 32

11.5. Paso al límite inductivo en hipercohomología \dotfill 35

11.6. Hiperhomología de un funtor respecto de un complejo $K_\bullet$
\dotfill 39

11.7. Hiperhomología de un funtor respecto de un bicomplejo
$K_{\bullet\bullet}$ \dotfill 41

11.8. Suplemento sobre la cohomología de los complejos simpliciales \dotfill 43

11.9. Un lema sobre complejos de tipo finito \dotfill 46

11.10. Característica de Euler--Poincaré de un complejo de módulos de longitud finita
\dotfill 48

\medskip
\S~12. Suplemento sobre la cohomología de haces \dotfill 49

12.1. Cohomología de haces de módulos sobre espacios anillados \dotfill 49

12.2. Imágenes directas superiores \dotfill 57

12.3. Suplementos sobre los funtores Ext de haces \dotfill 60

12.4. Hipercohomología del funtor de imagen directa \dotfill 62

\oldpage[III]{166}

\S~13. Límites proyectivos en álgebra homológica \dotfill 64

13.1. La condición de Mittag--Leffler \dotfill 64

13.2. La condición de Mittag--Leffler para grupos abelianos \dotfill 65

13.3. Aplicación: cohomología de un límite proyectivo de haces \dotfill 68

13.4. Condición de Mittag--Leffler y objetos graduados asociados a sistemas
proyectivos \dotfill 69

13.5. Límites proyectivos de sucesiones espectrales de complejos filtrados
\dotfill 71

13.6. Sucesión espectral de un funtor relativa a un objeto dotado de una
filtración finita \dotfill 73

13.7. Funtores derivados de un límite proyectivo de argumentos \dotfill 75

\medskip
\textsc{Capítulo III.} --- \textbf{Estudio cohomológico de los haces coherentes}
\dotfill 81

\medskip
\S~1. Cohomología de los esquemas afines \dotfill 82

1.1. Recordatorio sobre el complejo de álgebra exterior \dotfill 82

1.2. Cohomología de \v Cech de un recubrimiento abierto \dotfill 85

1.3. Cohomología de un esquema afín \dotfill 88

1.4. Aplicación a la cohomología de preesquemas cualesquiera \dotfill 89

\medskip
\S~2. Estudio cohomológico de los morfismos proyectivos \dotfill 95

2.1. Cálculos explícitos de ciertos grupos de cohomología \dotfill 95

2.2. El teorema fundamental de los morfismos proyectivos \dotfill 100

2.3. Aplicación a haces graduados de álgebras y módulos \dotfill 102

2.4. Una generalización del teorema fundamental \dotfill 107

2.5. Característica de Euler--Poincaré y polinomio de Hilbert \dotfill 109

2.6. Aplicación: criterios de amplitud \dotfill 111

\medskip
\S~3. Teorema de finitud para los morfismos propios \dotfill 115

3.1. El lema de dévissage \dotfill 115

3.2. El teorema de finitud: caso de los esquemas ordinarios \dotfill 116

3.3. Generalización del teorema de finitud (esquemas ordinarios) \dotfill 118

3.4. Teorema de finitud: caso de los esquemas formales \dotfill 119

\medskip
\S~4. El teorema fundamental de los morfismos propios. Aplicaciones \dotfill 122

4.1. El teorema fundamental \dotfill 122

4.2. Casos particulares y variantes \dotfill 129

4.3. Teorema de conexión de Zariski \dotfill 130

4.4. «Teorema principal» de Zariski \dotfill 135

4.5. Compleciones de módulos de homomorfismos \dotfill 138

4.6. Relaciones entre morfismos formales y morfismos ordinarios \dotfill 139

4.7. Un criterio de amplitud \dotfill 145

4.8. Morfismos finitos de preesquemas formales \dotfill 146

\oldpage[III]{167}

\S~5. Un teorema de existencia para los haces algebraicos coherentes \dotfill 149

5.1. Enunciado del teorema \dotfill 149

5.2. Demostración del teorema de existencia: caso proyectivo y cuasiproyectivo
\dotfill 151

5.3. Demostración del teorema de existencia: caso general \dotfill 154

5.4. Aplicación: comparación entre morfismos de preesquemas ordinarios y
morfismos de preesquemas formales. Preesquemas formales algebraizables \dotfill 156

5.5. Una descomposición de ciertos preesquemas \dotfill 159

\medskip
\textsc{Bibliografía (continuación)} \dotfill 161

\textsc{Índice de notaciones} \dotfill 162

\textsc{Índice terminológico} \dotfill 163

\begin{flushright}
\emph{Recibido el 28 de junio de 1961.}
\end{flushright}
\end{flushleft}
